%% LyX 2.4.4 created this file.  For more info, see https://www.lyx.org/.
%% Do not edit unless you really know what you are doing.
\documentclass[english]{article}
\usepackage[T1]{fontenc}
\usepackage[utf8]{inputenc}
\setcounter{secnumdepth}{-2}
\setlength{\parindent}{0bp}
\usepackage{amsmath}
\usepackage{amsthm}
\usepackage{amssymb}

\makeatletter
%%%%%%%%%%%%%%%%%%%%%%%%%%%%%% Textclass specific LaTeX commands.
\theoremstyle{definition}
\newtheorem*{problem*}{\protect\problemname}

\@ifundefined{date}{}{\date{}}
%%%%%%%%%%%%%%%%%%%%%%%%%%%%%% User specified LaTeX commands.
\usepackage{graphicx}
\usepackage{geometry}
\newcommand{\limi}{\underset{n\rightarrow\infty}{\lim}}
\newcommand{\limxz}{\underset{x\rightarrow x_{0}}{\lim}}
\newcommand{\limfxz}{\underset{x\rightarrow x_{0}}{\lim}f(x)}
\newcommand{\limfxm}{\underset{x\rightarrow x_{0}^{-}}{\lim}f(x)}
\newcommand{\limfxp}{\underset{x\rightarrow x_{0}^{+}}{\lim}f(x)}
\newcommand{\limxm}{\underset{x\rightarrow x_{0}^{-}}{\lim}}
\newcommand{\limxp}{\underset{x\rightarrow x_{0}^{+}}{\lim}}
\newcommand{\sqrm}{M_{m\times m}(\mathbb{F})}
\newcommand{\seqa}{(a_{n})_{n=1}^{\infty}}
\newcommand{\seqx}{(x_{n})_{n=1}^{\infty}}
\newcommand{\funcdr}{f:D\rightarrow\mathbb{R}}
\newcommand{\der}{\limxz\frac{f(x)-f(x_{0})}{x-x_{0}}}
\usepackage[all]{pdfprivacy}

\makeatother

\usepackage{babel}
\providecommand{\problemname}{Problem}

\begin{document}
\title{{\Large HUJI, Infinitesimal Calculus 1, Semester 2026A, Exam A, Q1\vspace{-2em}}}
\maketitle
\begin{problem*}
Prove that a bounded sequence is converging if and only if it has
exactly one partial limit. 
\end{problem*}
\begin{proof}
We prove the two implications separately.

\medskip{}

\textbf{First direction:} 

Let $\left(a_{n}\right)_{n=1}^{\infty}$ be a bounded sequence and
let us assume that it's converging to $L\in\mathbb{R}$.

\emph{Existence:} by the Bolzano-Weierstrass theorem, $\left(a_{n}\right)$
has a converging sub-sequence.

\emph{Uniqueness:} Let $\left(a_{n_{k}}\right)_{k=1}^{\infty}$ be
a converging sub-sequence of $\left(a_{n}\right).$

Then from the inheritance theorem (every sub-sequence of a convergent
sequence converges to the same limit), it also converges to $L$,
as needed.

\medskip{}

\textbf{Second direction:}

Let $\left(a_{n}\right)_{n=1}^{\infty}$ be a bounded sequence such
that it has exactly one partial limit, $L\in\mathbb{R}$.

Let us assume for the sake of contradiction that $\left(a_{n}\right)_{n=1}^{\infty}$
doesn't converge.

In particular, it doesn't converge to $L$, so there exists $\varepsilon_{1}>0$
such that $|a_{n}-L|>\varepsilon_{1}$ is a property that holds for
infinitely many $n$ (otherwise it would converge to $L$).

Since it's a frequent property of $\left(a_{n}\right)$, there exists
a sub-sequence $\left(a_{n_{j}}\right)_{j=1}^{\infty}$ of which all
elements fulfill this property.

In particular, $\left(a_{n_{j}}\right)_{j=1}^{\infty}$ inherits the
boundedness, therefore by the Bolzano-Weierstrass theorem it has a
converging sub-sequence $\left(a_{n_{j_{w}}}\right)_{w=1}^{\infty}$
, i.e. there exists $L'\in\mathbb{R}$ such that $\underset{w\rightarrow\infty}{\lim}a_{n_{j_{w}}}=L'$.

However, every element of this sub-sequence fulfills the property
of being far from $L$ by at least $\varepsilon_{1}$, and since limits
weakly preserve sequence inequalities, $|L'-L|\geq\varepsilon_{1}>0$,
therefore $L'\neq L$.

$L'$ is also a partial limit of $\left(a_{n}\right)$, since the
sub-sequence of a sub-sequence is also a sub-sequence, which contradicts
our assumption.
\end{proof}

\end{document}
